(1) 每段弦的小振幅方程 \(\mu_j y_{tt}=Ty_{xx}\)，故 \(c_j=\sqrt{T/\mu_j}\)，\(k_j=\omega/c_j\)。沿正向的机械功率与速度幅值满足 \(\overline P_j=\frac12Z_j|v_y|^2\)，其中 \(v_y=y_t\)、特性机械阻抗 \(Z_j=T/c_j=\sqrt{T\mu_j}\)。同一频率下位移幅度 \(A\) 的时间平均功率为 \(\overline P_j=\frac12 Z_j\omega^2|A|^2\)，这为界面能量检验提供了正确权重，不可直接用位移幅度平方比较不同密度的功率。

(2) 写左侧 \(y_1=A_ie^{i(k_1x-\omega t)}+A_re^{i(-k_1x-\omega t)}\)，右侧 \(y_2=A_te^{i(k_2x-\omega t)}\)。界面位移连续给 \(A_i+A_r=A_t\)，横向张力 \(T y_x\) 连续给 \(k_1(A_i-A_r)=k_2A_t\)。故反射系数 \(r=(k_1-k_2)/(k_1+k_2)=(Z_1-Z_2)/(Z_1+Z_2)\)，透射振幅 \(t=2Z_1/(Z_1+Z_2)\)。反射率 \(R=|r|^2\)，透射率 \(\mathcal T=(Z_2/Z_1)|t|^2=4Z_1Z_2/(Z_1+Z_2)^2\)，两者相加正好一。

(3) 对于匹配层，在其左端将向右与向左行波叠加，利用左右两端的位移和横向力连续，得到入射端等效阻抗 \(Z_{\rm in}=Z_m[Z_2+iZ_m\tan(k_md)]/[Z_m+iZ_2\tan(k_md)]\)（取时间相位 \(e^{-i\omega t}\) 时可整体改变虚部符号，不改变四分之一波长结果）。当 \(k_md=\pi/2\) 模 \(\pi\) 时，极限 \(Z_{\rm in}=Z_m^2/Z_2\)。令此值等于 \(Z_1\) 消除首端反射，得到 \(Z_m=\sqrt{Z_1Z_2}\)，对应统一拉力下 \(\mu_m=\sqrt{\mu_1\mu_2}\)。厚度 \(d=(2n+1)\lambda_m/4\) 给出匹配频率。

(4) 匹配层通过相位相消而非消除所有界面的单独反射，所以当频率改变，\(k_md\ne\pi/2\) 时 \(Z_{\rm in}\ne Z_1\)，反射一般恢复。若左右阻抗本来相同，即 \(\mu_1=\mu_2\)，单界面已全透射，不需要额外层；取 \(\mu_m=\mu_1\) 时任意厚度均不会引入反射。阻抗 \(Z_j=\sqrt{T\mu_j}\) 的量纲为千克每秒；几何平均的阻抗仍有相同量纲，四分之一波长厚度由 \(d=c_m\pi/(2\omega)\) 给出长度量纲。
